% =====================================================================
%  Partie 5 (~3 min) : l'IA comme outil -- décoder le cerveau
% =====================================================================
\partie{L'IA comme outil}{L'IA comme outil : décoder le cerveau}

% ---------------------------------------------------------------------
\begin{frame}{Redonner la parole}
  \begin{columns}[c,onlytextwidth]
    \column{.50\textwidth}
    \centering
    \begin{tikzpicture}[
        bx/.style={rounded corners=4pt, fill=bordeaux, text=white, align=center, minimum height=1.0cm, text width=2.4cm, font=\small},
        ob/.style={rounded corners=4pt, fill=fondclair, draw=jaune, text=jaune, align=center, minimum height=0.6cm, minimum width=2.0cm, font=\scriptsize},
        ar/.style={-{Stealth}, jaune, line width=1pt}]
      \node[bx] (c) at (0,0.2) {Cortex moteur\\{\scriptsize 253 électrodes}};
      \node[bx, fill=orange!50!fondclair] (n) at (0,-1.6) {Réseau de neurones};
      \node[ob] (t) at (-2.3,-3.0) {texte};
      \node[ob] (v) at (0,-3.0) {voix};
      \node[ob] (a) at (2.3,-3.0) {avatar};
      \draw[ar] (c)--(n); \draw[ar] (n)--(t); \draw[ar] (n)--(v); \draw[ar] (n)--(a);
    \end{tikzpicture}
    \column{.48\textwidth}
    \small
    \begin{itemize}
      \item<1-> Une patiente paralysée \textbf{tente} de parler, sans émettre de son.
      \item<2-> Le décodeur reconnaît des \textbf{phonèmes}, puis un modèle de langage choisit les phrases les plus probables.
      \item<3-> On retrouve les deux ingrédients de cet exposé : réseaux de neurones et \textbf{prédiction du mot suivant}.
    \end{itemize}
  \end{columns}
  \vspace{0.5em}
  \visible<4->{%
  \hspace*{3.6cm}\begin{tikzpicture}[x=0.042cm,
      lab/.style={anchor=east, font=\scriptsize, text=texte},
      val/.style={anchor=west, font=\scriptsize, text=gris}]
    \foreach \l/\v/\y/\c in {{Avant (sélection de lettres)}/14/0/gris, {Willett et al. 2023}/62/-0.42/orange, {Metzger et al. 2023}/78/-0.84/orange, {Conversation naturelle}/160/-1.26/jaune}{
      \node[lab] at (-2,\y) {\l};
      \fill[\c] (0,\y-0.15) rectangle (\v,\y+0.15);
      \node[val] at (\v+2,\y) {\v};}
    \node[font=\scriptsize, text=gris, anchor=west] at (95,-0.2) {mots par minute};
  \end{tikzpicture}}
  \src{Metzger et al., \emph{Nature} 2023 ; Willett et al., \emph{Nature} 2023}
\end{frame}

% ---------------------------------------------------------------------
\begin{frame}{Décoder sans chirurgie ?}
  \begin{columns}[T,onlytextwidth]
    \column{.48\textwidth}
    \carte[bleu]{IRMf : le sens d'une histoire}{%
      Un décodeur retrouve \textbf{l'idée générale} d'une histoire entendue ou imaginée, pas les mots exacts.\\[0.3em]
      {\color{gris}16 h d'enregistrement par personne ; ne marche que si la personne coopère.}\\[0.2em]
      {\scriptsize\color{gris}Tang et al., \emph{Nat. Neurosci.} 2023}}
    \column{.48\textwidth}
    \visible<2->{\carte[bleu]{MEG / EEG : la frappe au clavier}{%
      Décodage des phrases tapées par 35 volontaires.\\[0.3em]
      {\color{gris}Erreur sur les caractères : 32\,\% en MEG, 67\,\% en EEG.}\\[0.2em]
      {\scriptsize\color{gris}Lévy et al. (Meta), 2025}}}
  \end{columns}
  \vspace{0.6em}
  \visible<3->{\small\textbf{Questions éthiques} : consentement, \textbf{vie privée mentale}, accès équitable à ces technologies.}
  \vspace{0.3em}
  \visible<4->{\cle{Des résultats réels, mais encore loin de « lire dans les pensées ».}}
\end{frame}
